\rok{Новембар 2024.}{B}

Први практични тест из Алгоритама и структура података

\zadatak{1}{Insertion Sort}
Написати \textit{Insertion Sort} алгоритам за сортирање целобројног низа дужине $n$.

\textbf{Додатне напомене за решавање задатка:}
\begin{itemize}
    \item Улаз је несортиран низ целих бројева.
    \item Излаз је сортиран низ целих бројева.
    \item У шаблону се налази класа \texttt{RandomArrayGenerator} коју треба искористити за генерисање низа случајних бројева.
    \item Креирати класу \texttt{InsertionSort} и у оквиру ње генерисати методу:
    \begin{Verbatim}[fontsize=\footnotesize, numbers=left, xleftmargin=10mm]
public static void Sort(long[] arr)
    \end{Verbatim}
    \item Обезбедити код у \texttt{Main} методи за тестирање алгоритма.
\end{itemize}



\zadatak{2}{Спајање две једноструко спрегнуте листе, опадајући поредак}
Спојити две претходно сортиране, једноструко спрегнуте листе у једну, тако да елементи резултујуће листе буду сортирани у опадајућем поретку. Алгоритам је неопходно написати са линеарном временском комплексношћу $O(n)$.

\textbf{Додатни захтеви за решавање задатка:}
\begin{itemize}
    \item Алгоритам писати у оквиру класе \texttt{LinkedList}.
    \item Претпоставимо да су почетне листе већ сортиране.
    \item Захтеван потпис методе:
    \begin{Verbatim}[fontsize=\footnotesize, numbers=left, xleftmargin=10mm]
public static LinkedList MergingTwoLinkedLists(LinkedList list1,
                                     LinkedList list2)
    \end{Verbatim}
\end{itemize}

\textbf{Примери:}
\begin{itemize}
    \item \textbf{Улаз:} $4 \leftrightarrow 3 \leftrightarrow 2 \leftrightarrow 1$ ; \quad $4 \leftrightarrow 3 \leftrightarrow 3 \leftrightarrow 2$ \\
          \textbf{Излаз:} $4 \leftrightarrow 4 \leftrightarrow 3 \leftrightarrow 3 \leftrightarrow 3 \leftrightarrow 2 \leftrightarrow 2 \leftrightarrow 1$
    \item \textbf{Улаз:} $7 \leftrightarrow 6 \leftrightarrow 5 \leftrightarrow 4$ ; \quad $9 \leftrightarrow 8 \leftrightarrow 4 \leftrightarrow 1$ \\
          \textbf{Излаз:} $9 \leftrightarrow 8 \leftrightarrow 7 \leftrightarrow 6 \leftrightarrow 5 \leftrightarrow 4 \leftrightarrow 4 \leftrightarrow 1$
\end{itemize}



\zadatak{3}{Производ бројева реда једнак $k$}
Написати алгоритам који проналази све узастопне поднизове елемената у реду чији је производ једнак задатом броју $k$. Метода треба да испише сваки подниз који задовољава овај услов.

\textbf{Додатни захтеви за решавање задатка:}
\begin{itemize}
    \item Задатак решавати помоћу структуре реда (класа \texttt{Queue}) која је дата у оквиру шаблона задатка.
    \item Захтеван потпис методе:
    \begin{Verbatim}[fontsize=\footnotesize, numbers=left, xleftmargin=10mm]
public void FindSubsequencesWithProduct(int targetProduct)
    \end{Verbatim}
\end{itemize}

\textbf{Пример 1:}
\begin{itemize}
    \item \textbf{Улаз за ред:} 2, 12, 6, 8, 3 ; $k = 24$
    \item \textbf{Испис:}
    \begin{itemize}
        \item Подниз производа 24: 2, 12
        \item Подниз производа 24: 8, 3
    \end{itemize}
\end{itemize}

\textbf{Пример 2:}
\begin{itemize}
    \item \textbf{Улаз за ред:} 9, 6, 6, 1, 3 ; $k = 36$
    \item \textbf{Испис:}
    \begin{itemize}
        \item Подниз производа 36: 6, 6
    \end{itemize}
\end{itemize}

\sablon{AISP2024_Test1_GrupaB}{2024/Novembar/GrupaB/AISP2024_Test1_GrupaB.linq}
